\section{Estimating labour from a sample}
\label{sec:sampling}

Divide each contributor's observation period into disjoint intervals. Index the
eligible person-intervals by $U=\{1,\ldots,M\}$. Interval $m$ lasts $d_m$ hours
and contains $y_m\in[0,d_m]$ person-hours of work on the requirements. The target
is $H=\sum_{m\in U}y_m$. Select a sample $s$ with known inclusion probability
$\pi_m>0$ for every interval. Draw after the work period if asking during it could interrupt the activity. Inference holds the completed work history fixed and
averages over possible draws.

\subsection{Use predictions to reduce sampling variation}

Fix a prediction $f_m\in[0,d_m]$ for each interval before the draw and write
$r_m=y_m-f_m$ and $F=\sum_{m\in U}f_m$. If the selected hours are known, estimate
the total by
\begin{equation}
  \widehat H=F+\sum_{m\in s}\frac{r_m}{\pi_m}.
  \label{eq:short-difference}
\end{equation}
The indicator that interval $m$ is selected has expectation $\pi_m$. Thus the
expected weighted correction is $\sum_{m\in U}r_m$, and
$\mathbb{E}[\widehat H]=H$. Selections may depend on one another. Predictions
may have systematic errors because the correction estimates their total error.
Setting each $f_m$ to zero yields the Horvitz--Thompson estimator \cite{horvitz1952}.
The generalised difference estimator can reduce variation across draws
when predictions leave a less variable correction \cite{cassel1976}.

For a simple random sample drawn without replacement, suppose
$1\le n\le M$ and $M>1$. Then the variance of \eqref{eq:short-difference} is
$K_n S_r^2$, where $K_n=M^2(1-n/M)/n$ and $S_r^2$ is the population variance of
the residuals with divisor $M-1$. This follows from $\pi_m=n/M$ and the
finite-population variance of the sample mean. The direct estimator has
variance $K_n S_y^2$, with the hours themselves in place of the residuals.
For $n<M$, prediction correction improves precision exactly when
$S_r^2<S_y^2$. A census has zero sampling variance.

\subsection{Follow up missed observations or bound their effect}

Let $A$ be the intervals with usable initial answers. Randomly select $B$
from $s\setminus A$ so that each missed interval has known conditional
selection probability $\eta_m>0$. If both stages supply accurate hours for the
intervals they cover, the estimate is
\[
  F+\sum_{m\in A}\frac{r_m}{\pi_m}
   +\sum_{m\in B}\frac{r_m}{\pi_m\eta_m}.
\]
Conditional on $s$ and $A$, the second draw recovers the complete-sample
correction in \eqref{eq:short-difference} in expectation. Taking the expectation
over the first draw recovers $H$.
Workload may affect the initial answer rate \cite{nonresponsefollowup}.

When accurate hours remain unavailable, bound each selected interval by $[l_m,u_m]$ from the available observations. For a missed answer, the
physical limits $[0,d_m]$ suffice. The weighted estimate increases with each $y_m$. Evaluating it at the lower
limits yields $\widehat H^-$ and at the upper limits yields $\widehat H^+$.
The estimate that would have used the true sampled hours lies between these
values whenever every supplied range contains its true hours. These bounds
account for incomplete measurement on the selected intervals.

For simple random sampling, add a finite-sample margin for the draw itself.
All residuals lie between $a=-\max_{m\in U} f_m$ and
$b=\max_{m\in U}(d_m-f_m)$. The sum of squared deviations from their mean
is at most the sum of squared deviations from $(a+b)/2$, so
$S_r^2\le M(b-a)^2/[4(M-1)]$. Chebyshev's inequality gives a margin
\[
  e=\sqrt{\frac{K_n M(b-a)^2}{4(M-1)\alpha}}\qquad (0<\alpha<1).
\]
The interval $[\widehat H^--e,\widehat H^++e]$, intersected with the physical range
$[\sum_{m\in s}l_m,\;\sum_{m\in s}u_m+\sum_{m\notin s}d_m]$, covers $H$ with probability at least $1-\alpha$ when the
measurement limits hold on every draw. Use $e=0$ for a census, including
$M=1$. If the chance that any sample range excludes its true value is
at most $\beta$, coverage is at least $1-\alpha-\beta$ by the union
bound. The margin can be wide when little is known. If eligible work
outside the list has a defensible total capacity $D$, adding $D$ to the upper
limit covers the expanded target at the same level. A separate estimate
requires evidence about the time those contributors spent on the requirements.
